\subsection{Regularity of Cographic Matroids}

\begin{lemma}[Row space of a standard representation]
    \label{lem:row-space-standard}
    Let $X$ and $Y$ be disjoint finite sets and let
    
    \[
      B \in \mathbb{F}_2^{X \times Y}.
    \]

    Consider the matrix
    
    \[ A := [\, \mathbf{1}_x \mid B \,] \in \mathbb{F}_2^{X \times (X \cup Y)},\]
    
    where the columns are indexed by $E := X \cup Y$ and the rows by $X$.
    Then the row space of $A$ is
    
    \[ \text{row}(A) \;=\; \{\, (u,\, uB) \mid u \in \mathbb{F}_2^X \,\} \;\subseteq\; \mathbb{F}_2^X \oplus \mathbb{F}_2^Y \cong \mathbb{F}_2^{E}. \]
\end{lemma}

\begin{proof}
    The $x$-th row of $A$ is $(e_x, B_{x,*})$, where $e_x$ is the standard basis vector
    in $\mathbb{F}_2^X$ and $B_{x,*}$ is the $x$-th row of $B$.
    A general linear combination of the rows is therefore
    \[
      \sum_{x \in X} u_x (e_x, B_{x,*})
      = \bigl( u,\, \sum_{x \in X} u_x B_{x,*} \bigr)
      = (u,\, uB),
    \]
    where $u = (u_x)_{x \in X} \in \mathbb{F}_2^X$.
    Conversely, every pair $(u, uB)$ arises in this way, so these are exactly the row vectors.
\end{proof}

\begin{lemma}[Orthogonal complement of a standard row space]
    \label{lem:orth-complement-standard}
    Let $A = [\mathbf{1}_x \mid B]$ be as in Lemma~\ref{lem:row-space-standard}, and let
    \[U := \text{row}(A) \subseteq \mathbb{F}_2^{X \cup Y}.\]
    Then the orthogonal complement of $U$ is
    \[U^\perp = \{\, (b B^{\mathsf T},\, b) \mid b \in \mathbb{F}_2^Y \,\}.\]
    Equivalently, if $B^* := -B^{\mathsf T}$, then
    \[U^\perp= \{\, (b B^*,\, b) \mid b \in \mathbb{F}_2^Y \,\}.
    \]
\end{lemma}

\begin{proof}
    Write vectors in $\mathbb{F}_2^{X \cup Y}$ as pairs $(a,b)$ with
    $a \in \mathbb{F}_2^X$ and $b \in \mathbb{F}_2^Y$.
    By Lemma~\ref{lem:row-space-standard}, any element of $U$ has the form
    $(u, uB)$ with $u \in \mathbb{F}_2^X$.
    The orthogonality condition $(a,b) \in U^\perp$ means
    \begin{align*}
      0 &= (a,b) \cdot (u, uB) \\ 
        &= a \cdot u + b \cdot (uB) \\
        &= a \cdot u + (b B^{\mathsf T}) \cdot u \\
        &= (a + b B^{\mathsf T}) \cdot u
    \end{align*}
    for all $u \in \mathbb{F}_2^X$.
    Hence we must have $a = b B^{\mathsf T}$, and then
    \[
      U^\perp = \{\, (b B^{\mathsf T}, b) \mid b \in \mathbb{F}_2^Y \,\}.
    \]
    Over $\mathbb{F}_2$ we have $-1 = 1$, so $B^* = -B^{\mathsf T} = B^{\mathsf T}$,
    yielding the alternative description.
\end{proof}

\begin{lemma}[Row space of the dual standard matrix]
    \label{lem:row-space-dual-standard}
    With $B$ and $B^* = -B^{\mathsf T}$ as above, define
    \[A^* := [\, \mathbf{1}_y \mid B^* \,] \in \mathbb{F}_2^{Y \times (X \cup Y)}.\]
    Then
    \[\text{row}(A^*) = U^\perp,\]
    where $U = \text{row}(A)$ and $U^\perp$ is given by
    Lemma~\ref{lem:orth-complement-standard}.
\end{lemma}

\begin{proof}
    The $y$-th row of $A^*$ is $(e_y, B^*_{y,*})$ with $e_y \in \mathbb{F}_2^Y$.
    A general linear combination of the rows is
    \[\sum_{y \in Y} b_y (e_y, B^*_{y,*}) = \bigl( b,\, b B^* \bigr),
    \]
    where $b = (b_y)_{y \in Y} \in \mathbb{F}_2^Y$.
    Thus
    \[\text{row}(A^*) = \{\, (b, b B^*) \mid b \in \mathbb{F}_2^Y \,\}.\]
    Identifying $\mathbb{F}_2^{X \cup Y}$ as $\mathbb{F}_2^X \oplus \mathbb{F}_2^Y$
    with coordinates ordered as $(X, Y)$, this is exactly the set
    \[\{\, (b B^*, b) \mid b \in \mathbb{F}_2^Y \,\},\]
    which coincides with $U^\perp$ by Lemma~\ref{lem:orth-complement-standard}.
\end{proof}

\begin{lemma}[Dual vector matroid via orthogonal complement]
    \label{lem:dual-vectormatroid-orth}
    Let $A$ and $A'$ be matrices over a field $F$ with the same column index set $E$,
    and suppose
    \[\text{row}(A') = \text{row}(A)^{\perp} \subseteq F^{E}.\]
    Let $M(A)$ and $M(A')$ be the vector matroids represented by $A$ and $A'$. Then
    \[M(A') = M(A)^{*}.\]
\end{lemma}


\begin{proof}
    Let $F \subseteq E$.  

    \emph{($\Rightarrow$)}  
    Suppose $F$ is dependent in $M(A)$.  
    Then there exists a nonzero vector $c \in F^{F}$ such that $A_F c = 0$.  
    Extend $c$ by zero outside $F$ (still denoted $c$).  
    The condition $A c = 0$ means each row $r$ of $A$ satisfies $r \cdot c = 0$,  
    hence $c \in \text{row}(A)^{\perp} = \text{row}(A')$.  
    Write
    \[
        c = \sum_i \lambda_i r'_i,
    \]
    where the $r'_i$ are rows of $A'$ and not all $\lambda_i$ are zero.  
    For every $e \in E\setminus F$ we have $c_e = 0$, so
    \[
        \left(\sum_i \lambda_i r'_i\right)\big|_{E\setminus F} = 0.
    \]
    Hence the rows of $A'$ indexed by $E \setminus F$ admit a nontrivial
    linear combination giving the zero row, so $E \setminus F$ is dependent in $M(A')$.

    \emph{($\Leftarrow$)}  
    The same argument with $A$ and $A'$ interchanged, using
    $\text{row}(A) = (\text{row}(A')^\perp)$, shows that if
    $E\setminus F$ is dependent in $M(A')$, then $F$ is dependent in $M(A)$.

    Thus
    \[
        F \text{ dependent in } M(A)
        \;\Longleftrightarrow\;
        E \setminus F \text{ dependent in } M(A'),
    \]
    which is the defining property of duality.
\end{proof}





\begin{theorem}[Dual of standard representation corresponds to dual matroid]
\label{thm:StandardRepr.toMatroid_dual}
Let M be a binary matroid on ground set $E = X \cup Y$, with
standard representation $B$ so that
\[A = [\, \mathbf{1}_X \mid B \,].\]
Let $B^{*} := -B^{\mathsf T}$ and
\[A^{*} := [\, \mathbf{1}_Y \mid B^{*} \,].\]
Then $M(A^{*}) = M(A)^{*} = M^{*}$.
\end{theorem}

\begin{proof}
    By Lemma~\ref{lem:row-space-standard} and Lemma~\ref{lem:orth-complement-standard},
    if $U = \text{row}(A)$ then $U^\perp$ has the form
    \[U^\perp = \{\, (b B^*, b) \mid b \in \mathbb{F}_2^Y \,\}.\]
    By Lemma~\ref{lem:row-space-dual-standard}, we have
    \[\text{row}(A^*) = U^\perp = \text{row}(A)^\perp.\]
    Therefore, by Lemma~\ref{lem:dual-vectormatroid-orth},
    the column-matroid $M(A^*)$ is the dual of $M(A)$:
    \[M(A^*) = M(A)^* = M^*.\]
\end{proof}

\begin{lemma}
    \label{lem:Matroid.Dual.isRegular}
    The dual matroid of a regular matroid is also a regular matroid.
\end{lemma}

\begin{proof}
    Let $M$ be a regular matroid. We wish to show that $M^{*}$ is also regular. 
    
    Take a standard $ \mathbb Z_2$-representation matrix $B$ of $M$. By Lemma~\ref{StandardRepr.toMatroid_isRegular_iff_hasTuSigning}, since $M$ is regular, there exists a TU signing $B'$ of $B$: $B'$ is a matrix over $\mathbb Q$ that is TU, and $|B'(i,j)| = B(i,j)$ for all entries.  So $M$ is represented (over $\mathbb Q$) by a TU matrix $B'$ whose pattern of zero and non-zero entries is exactly that of $B$.

    From Theorem \ref{thm:StandardRepr.toMatroid_dual}, if a matroid $M$ has standard representation matrix $B$, then its dual $M^{*}$ has the standard representation matrix $B^{*} = -B^{\intercal}$. The TU signing of this dual standard matrix, $(B')^{*} = -(B')^{\intercal}$, preserves total unimodularity, so $(B')^{*}$ is a TU matrix whose support is exactly $B^{*}$.

    Since we have just exhibited a TU signing of $M^{*}$ (i.e., $(B')^{*}$), the dual matroid $M^{*}$ is regular by Lemma~\ref{StandardRepr.toMatroid_isRegular_iff_hasTuSigning}.
\end{proof}

\begin{theorem}
    \label{Matroid.IsCographic.isRegular}
    Every cographic matroid is regular.
\end{theorem}

\begin{proof}
    We know that all graphic matroids are regular by Theorem~\ref{Matroid.IsGraphic.isRegular}. Recall that we say a matroid is cographic if its dual is graphic. So it suffices to show regularity is preserved under duals, which we showed in Lemma \ref{lem:Matroid.Dual.isRegular}.
\end{proof}